\clearpage
\suppressfloats[t]
\appendixsection{\revisionr3{Additional Results for the Attention and Norm Comparison}}
\label{app:paired_diagnostics}
\revisionr3{Across the 16,000 selected records, we retain 252 generations reaching the 256-token cap, nine empty evaluated answers, and seven no-output-state zero fallbacks (all seven in Gemma/NQ). These sets can overlap. No record is excluded from the paired metrics.}

\revisionr3{These analyses use the same fixed first-answer-line samples as Table~\ref{tab:paired_s}, without score filtering or class balancing. Table~\ref{tab:paired_r} reports ROUGE-L-based metrics; Table~\ref{tab:paired_differences} reports paired AUC$_s$ differences. PCC in the detector tables uses continuous reference similarity, whereas the structural plots correlate detector scores with one another. These are distinct quantities.}
\begin{table*}[htbp]
\centering\small
\setlength{\tabcolsep}{4pt}
\revisionr3{\textbf{Panel A: HIDE full-benchmark reference (Appendix B)}\par\smallskip}
\revisionr3{\begin{tabular}{llrrr}
\toprule
Model & Dataset & $N$ & AUC$_r$ & PCC$_r$ \\
\midrule
Llama-3-8B & SQuAD & 5928 & 77.45 & 50.43 \\
Llama-3-8B & RACE & 3498 & 66.95 & 22.87 \\
Llama-3-8B & NQ & 3610 & 79.17 & 33.91 \\
Llama-3-8B & TriviaQA & 9960 & 66.82 & 35.19 \\
Gemma-2-9B & SQuAD & 5928 & 76.66 & 48.43 \\
Gemma-2-9B & RACE & 3498 & 57.99 & 16.50 \\
Gemma-2-9B & NQ & 3610 & 75.96 & 15.41 \\
Gemma-2-9B & TriviaQA & 9960 & 81.02 & 43.88 \\
\bottomrule\end{tabular}}
\par\medskip
\revisionr3{\textbf{Panel B: Matched first-answer-line comparison}\par\smallskip}
\revisionr3{\begin{tabular}{llr rr rr rr}
\toprule
& & & \multicolumn{2}{c}{HIDE} & \multicolumn{2}{c}{Attention $\Omega$} & \multicolumn{2}{c}{Norm $\widetilde\Delta_{\rm in}$} \\
Model & Dataset & $N_+$ & AUC$_r$ & PCC$_r$ & AUC$_r$ & PCC$_r$ & AUC$_r$ & PCC$_r$ \\
\midrule
Llama-3-8B & SQuAD & 922 & 77.32 & 48.97 & 85.57 & 60.71 & 51.24 & -0.69 \\
Llama-3-8B & RACE & 875 & 58.08 & 9.35 & 59.88 & 14.16 & 47.00 & -11.84 \\
Llama-3-8B & NQ & 280 & 78.58 & 31.79 & 84.00 & 32.58 & 71.44 & 22.36 \\
Llama-3-8B & TriviaQA & 1269 & 58.58 & 5.56 & 73.38 & 43.02 & 64.69 & 26.77 \\
\midrule
Gemma-2-9B & SQuAD & 849 & 76.51 & 48.85 & 87.58 & 65.46 & 54.13 & 9.15 \\
Gemma-2-9B & RACE & 1185 & 58.12 & 16.29 & 50.41 & -1.90 & 50.95 & 1.93 \\
Gemma-2-9B & NQ & 480 & 63.16 & 16.47 & 69.10 & 23.10 & 60.08 & 13.96 \\
Gemma-2-9B & TriviaQA & 1359 & 51.18 & -5.81 & 66.30 & 32.04 & 62.86 & 21.04 \\
\bottomrule
\end{tabular}}
\caption{\revisionr3{Panel A reproduces HIDE values from Table~\ref{tab:results_exp} of Appendix B without alteration ($N$: full-benchmark sample size). Panel B evaluates all three scores on the same 2,000 fixed seed-42 examples per model/dataset under first-answer-line stopping; $N_+$ refers only to Panel B. Because cohorts and stopping differ, Panel A is a reference, not a matched comparator for Panel B. $N_+$ counts answers with ROUGE-L greater than 0.5. AUC and PCC are multiplied by 100; PCC uses the continuous reference metric. Within Panel B, all scores use the same examples and labels; higher raw scores predict correctness. No direction is selected using test performance.}}
\label{tab:paired_r}
\end{table*}

\begin{table*}[t]
\centering\small
\begin{tabular}{llrrr}
\toprule
Model & Dataset & HIDE $-$ attention & 95\% CI & HIDE $-$ norm \\
\midrule
Llama-3-8B & SQuAD & -9.01 & [-11.43, -6.69] & +22.55 \\
Llama-3-8B & RACE & -3.00 & [-5.27, -0.82] & +11.69 \\
Llama-3-8B & NQ & -5.77 & [-9.66, -2.21] & +6.42 \\
Llama-3-8B & TriviaQA & -16.67 & [-19.96, -13.56] & -7.67 \\
Gemma-2-9B & SQuAD & -10.82 & [-13.37, -8.56] & +21.58 \\
Gemma-2-9B & RACE & +6.69 & [+3.70, +9.69] & +7.37 \\
Gemma-2-9B & NQ & -5.23 & [-8.69, -1.86] & +4.11 \\
Gemma-2-9B & TriviaQA & -17.11 & [-20.84, -13.51] & -11.70 \\
\bottomrule\end{tabular}
\caption{Paired AUC$_s$ differences in percentage points on the corrected 2,000-example cohorts. Intervals use 2,000 paired example bootstrap draws (seed 42), are pointwise, and are not adjusted for multiple comparisons.}
\label{tab:paired_differences}\end{table*}


\revisionr3{For actual selected count $n\geq1$, substituting all-ones Gram matrices into the adapted score gives $c(n)=(n-1)/n^2$ (and we define $c(0)=0$ for the no-state fallback). Table~\ref{tab:count_control} compares this diagnostic with HIDE on exactly the same examples. The near-unit correlations and small AUC differences indicate that the configured score is dominated by selected-token count in these runs; they do not establish a nonlinear-dependence advantage. We did not retune the kernel on the test labels.}
\begin{table*}[t]
\centering\small
\revisionr3{\begin{tabular}{llrrr}
\toprule
Model & Dataset & HIDE AUC$_s$ & Count-control AUC$_s$ & PCC(HIDE, control) \\
\midrule
Llama-3-8B & SQuAD & 76.33 & 76.38 & 1.000000 \\
Llama-3-8B & RACE & 61.66 & 61.92 & 1.000000 \\
Llama-3-8B & NQ & 78.47 & 78.57 & 1.000000 \\
Llama-3-8B & TriviaQA & 57.01 & 56.58 & 1.000000 \\
Gemma-2-9B & SQuAD & 75.90 & 75.92 & 0.998759 \\
Gemma-2-9B & RACE & 57.93 & 57.83 & 0.998718 \\
Gemma-2-9B & NQ & 62.74 & 63.01 & 0.999911 \\
Gemma-2-9B & TriviaQA & 49.83 & 49.78 & 0.999988 \\
\bottomrule\end{tabular}}
\caption{\revisionr3{Selected-token-count diagnostic on the same corrected cohorts. The control is $(n_{\rm eff}-1)/n_{\rm eff}^2$ for $n_{\rm eff}\geq1$, and zero otherwise. AUC is multiplied by 100; PCC is unscaled. This is a control for realized selected-token count, not raw output length.}}
\label{tab:count_control}\end{table*}

\begin{table*}[t]
\centering\small
\begin{tabular}{llrr}
\toprule
Model & Dataset & Mean input tokens & Mean generated tokens \\
\midrule
Llama-3.2-3B & SQuAD & 186.47 & 33.44 \\
Llama-3.2-3B & NQ & 20.30 & 17.34 \\
Llama-3-8B & SQuAD & 186.47 & 21.39 \\
Llama-3-8B & NQ & 20.30 & 20.50 \\
Gemma-2-9B & SQuAD & 186.20 & 39.29 \\
Gemma-2-9B & NQ & 21.14 & 22.16 \\
Gemma-2-27B & SQuAD & 186.20 & 45.65 \\
Gemma-2-27B & NQ & 21.14 & 12.14 \\
\bottomrule\end{tabular}
\caption{Mean token lengths for the timing queries in Table~\ref{tab:timing_revision}, under the earlier stopping protocol (maximum 256 generated tokens). Means include all queries and measured repeats; repeats share deterministic token sequences. These length differences limit a parameter-only interpretation of the scaling comparison.}
\label{tab:timing_lengths}\end{table*}

